\documentclass[11pt,a4paper]{article}

\usepackage[a4paper,margin=2.4cm]{geometry}
\usepackage[T1]{fontenc}
\usepackage{hyperref}
\usepackage{fancyhdr}

\hypersetup{
  colorlinks=true,
  linkcolor=blue,
  urlcolor=blue,
  pdftitle={EventHub Project Report},
  pdfauthor={EventHub Team}
}

\pagestyle{fancy}
\fancyhf{}
\fancyhead[L]{EventHub Project Report}
\fancyhead[R]{Draft Template}
\fancyfoot[C]{\thepage}
\setlength{\headheight}{14pt}

\setlength{\parindent}{0pt}
\setlength{\parskip}{0.65em}

\newcommand{\placeholder}[1]{\textit{#1}}

\title{\Huge EventHub Project Report}
\author{Team Members: \placeholder{Maksym DOLHOV and Mehdi AGHAEI}}
\date{\placeholder{final submission date --> 5 avril 2026!}}

\begin{document}

\maketitle

\section*{Purpose}

This document is a clean writing template for the final project report. Replace every placeholder sentence with the actual implementation details once the frontend, backend, integration, and deployment work are finished.

\tableofcontents
\newpage

\section{Introduction}

\placeholder{Summarize the project goal, the course context, the team organization, and the final expected deliverable.}

\section{Backend}

\subsection{Django}

\placeholder{Describe the Django backend architecture, the data model, the REST API, the business rules, the authentication and roles, and the testing strategy.}

\begin{itemize}
  \item \placeholder{Explain the Event, Participant, and Registration models.}
  \item \placeholder{Explain the dedicated many-to-many relation through Registration.}
  \item \placeholder{Explain validation rules and error handling.}
  \item \placeholder{Explain permissions, roles, and authentication once implemented.}
\end{itemize}

\subsection{Node.js and Express}

\placeholder{Describe the simplified comparative backend, its routes, its database choice, its middleware, and the analytical comparison with Django.}

\begin{itemize}
  \item \placeholder{State which entities were implemented in Express.}
  \item \placeholder{Explain the CRUD routes and middleware.}
  \item \placeholder{Compare Django and Node.js / Express in structure, philosophy, complexity, scalability, and ecosystem.}
\end{itemize}

\section{Frontend}

\subsection{HTML}

\placeholder{Describe the page structure, semantic layout, form markup, and how the main screens are organized.}

\subsection{CSS}

\placeholder{Describe the design system, layout rules, responsive behavior, and styling choices.}

\subsection{React}

\placeholder{Describe the component structure, hooks, routing, state handling, API calls, loading states, error states, token storage, and route protection.}

\section{Integration and Deployment}

\placeholder{Explain how the frontend connects to the Django backend, how environment variables are configured, how the production deployment works, and how the final public version was tested.}

\section{Conclusion}

\placeholder{Summarize what was built, what worked well, what was difficult, and what could be improved in a future iteration.}

\appendix

\section{Appendix}

\placeholder{Use the appendix for additional notes, commands, screenshots, comparisons, or supporting material that would interrupt the main flow of the report.}

\end{document}
